\section[Cote 63 : les quatre composantes de $j = 0$ ont quatre tangentes en $p = 2$]{Cote 63 : les quatre composantes de $j = 0$ ont quatre tangentes en $p = 2$\protect\verifie{Folder63}}\label{sec:63}

La cote 63, \emph{Formulaire courbes elliptiques}, rassemble des textes d'autres auteurs (la lettre de Tate à Cassels, le formulaire de Deligne d'après Tate, des tapuscrits d'Atkin) et des notes de Grothendieck ; la carte la date de 1967. Aux pages~95 à~99, il étudie le schéma $M_3$ des courbes elliptiques munies d'une base de leurs points d'ordre~$3$, au-dessus de $\mathbf Z[\frac13]$, et le lieu $j = 0$, c'est-à-dire $M'_3 = M_3 \cap V(c_4)$. Ce lieu a quatre composantes irréductibles, qui se rencontrent toutes en un point $a$ au-dessus de $\mathbf F_4$, en caractéristique~$2$. La page~97 pose la question : \q{À examiner si les 4 composantes de $(M'_3)_2$ ont \q{tangentes distinctes}} (\lot{63}{5}{97}). La page~99 commence par \q{\dots{} aux composantes irréductibles sont égales} et affirme que l'espace tangent de Zariski en $a$ se scinde en $V_1 \times V_2$, stable sous le groupe $G' = \mathrm{SL}(2, \mathbf F_3)$ (\lot{63}{5}{99}). La lecture en conclut, par un argument d'orbites, que les quatre branches ont la même tangente. Elles en ont quatre.

\noindent\emph{Le modèle.} On travaille sur la famille de Hesse $x^3 + y^3 + z^3 = 3t\,xyz$, de paramètre $t$, qui est le niveau~$3$ au-dessus de $\mathbf Z[\frac13, \zeta]$, $\zeta$ racine primitive cubique de l'unité. Pour cette famille, $c_4$ vaut $t(t^3 + 8)$ à une unité près (Artebani et Dolgachev) ; ce calcul et l'identification de la famille de Hesse à $M_3$ sont admis ici. Soient $R = \mathbf Z[\zeta]$, l'anneau $A = R[t]$, et l'idéal maximal $P = (2, t)$ du point $a$, de corps résiduel $A/P = R/2R = \mathbf F_4$. L'espace cotangent en $a$ est $P/P^2$, un $\mathbf F_4$-espace vectoriel de base les classes de $2$ et de~$t$. La tangente d'une branche lisse $f = 0$ passant par $a$ est la droite de $P/P^2$ engendrée par la classe de $f$.

\begin{encadre}
\begin{contrexemple}[Énoncé de l'édition, pages 97 à 99]
Dans $R[t]$, $t^3 + 8 = (t + 2)(t + 2\zeta)(t + 2\zeta^2)$, de sorte que $V(c_4)$ est formé des quatre branches $t = 0$, $t = -2$, $t = -2\zeta$, $t = -2\zeta^2$, toutes dans $P$. Pour deux branches distinctes $f$ et $g$ de cette liste et tous $a, b \in R[t]$, si $af + bg \in P^2$, alors $a \in P$ et $b \in P$ : les quatre tangentes sont deux à deux distinctes. L'involution de Hesse $t \mapsto (t + 2)/(t - 1)$ fixe $a$ et agit sur $P/P^2$ par $2 \mapsto 2$, $t \mapsto t + 2$ : c'est une transvection, qui ne respecte pas le scindage de la page~99.
\end{contrexemple}
\end{encadre}

\begin{proof}
La factorisation se vérifie en développant : le produit vaut $t^3 + 2(1 + \zeta + \zeta^2)t^2 + 4(\zeta + \zeta^2 + \zeta^3)t + 8\zeta^3$, et $1 + \zeta + \zeta^2 = 0$, $\zeta^3 = 1$.

Un polynôme $h = \sum h_i t^i$ de $R[t]$ est dans $P$ si et seulement si $2$ divise $h_0$ dans $R$ ; s'il est dans $P^2$, alors $4$ divise $h_0$ et $2$ divise $h_1$, car c'est vrai des produits de deux éléments de $P$ et ces conditions définissent un idéal. Soient $f = t + 2\alpha$ et $g = t + 2\beta$ deux branches, $\alpha \neq \beta$ dans $\{0, 1, \zeta, \zeta^2\}$, et $af + bg \in P^2$. Les coefficients de degré $0$ et $1$ de $af + bg$ sont $2(\alpha a_0 + \beta b_0)$ et $a_0 + b_0 + 2(\alpha a_1 + \beta b_1)$. Le premier est divisible par $4$ ; comme $R$ est un $\mathbf Z$-module libre de base $1, \zeta$, on peut simplifier par $2$, et $2$ divise $\alpha a_0 + \beta b_0$. Le second est divisible par $2$, donc $b_0 = 2w - a_0$ pour un $w \in R$, et $(\alpha - \beta) a_0 = 2(u - \beta w)$ si $\alpha a_0 + \beta b_0 = 2u$. Or chaque différence $\alpha - \beta$ est inversible modulo $2$ : $(-1)(-1) = 1$, $(-\zeta)(-\zeta^2) = 1$, $(1 - \zeta)(1 - \zeta^2) = 3 = 1 + 2$, et $(\zeta - \zeta^2)\,\zeta^2(1 - \zeta^2) = 3$. Si $(\alpha - \beta)\varepsilon = 1 + 2m$, on obtient $a_0 = 2\bigl(\varepsilon(u - \beta w) - m a_0\bigr)$, donc $a \in P$, puis $b_0 = 2w - a_0$ est pair et $b \in P$.

Pour l'involution, $t - 1$ n'est pas dans $P$ (son terme constant $-1$ n'est pas pair, en comparant les coordonnées sur la base $1, \zeta$), donc inversible dans l'anneau local en $a$, et $(t + 2) - (t - 1)(t + 2) = 4 - t^2 \in P^2$ : modulo $P^2$, l'image $(t + 2)/(t - 1)$ de $t$ est $t + 2$. Le nombre $2$ est fixe.
\end{proof}

\begin{remarque}
L'argument de la page~99 suppose que $G'$ respecte le scindage $V_1 \times V_2$ de l'espace tangent, $V_1$ vertical et $V_2$ horizontal, et qu'il agit trivialement sur $V_2$. L'involution de Hesse, qui provient du groupe $Q_8 \subset G'$, envoie la direction horizontale $t$ sur $t + 2$ : elle ne respecte pas le scindage, et les quatre droites non verticales de $P/P^2$ forment une seule orbite de quatre éléments au lieu d'être fixées. La multiplicité~$4$ du point et la non-régularité de $M'$ en $p = 2$ (page~95) ne sont pas en cause.
\end{remarque}

Ne sont pas démontrés ici : l'identification de la famille de Hesse au schéma $M_3$ et la formule de $c_4$, admises ; la structure de $M_3[\frac16] \cap V(c_4)$ de la page~98 ; les notes des pages~65 à~76 sur les périodes et les deux décompositions de $H^1_{\mathrm{dR}}$.
